\section{Discussion}

Webis-SR4ALL-26 provides a large-scale, cross-disciplinary corpus of systematic reviews with a uniform, fully open metadata backbone. By anchoring reviews and cited works to OpenAlex and by extracting method-related fields from available full texts, the corpus combines bibliographic context with structured descriptions of review objectives, eligibility criteria, and reported search strategies. This design prioritizes verifiability and cross-domain comparability, making it possible to study systematic review workflows and practices under consistent conditions across disciplines, while accepting uneven coverage due to incomplete full-text availability and heterogeneous reporting.
% The corpus operationalizes systematic review retrieval within a fully open and uniform infrastructure. This design prioritizes verifiability and cross-disciplinary comparability over database-specific fidelity, enabling controlled experimentation across disciplines under consistent conditions. By anchoring all records to OpenAlex metadata and indexed reference lists, the corpus provides a unified citation backbone while accepting variability in structured coverage due to incomplete full-text availability and heterogeneous reporting practices.

To illustrate one concrete use case, Section~\ref{sec:query-norm-demo} reported retrieval signals in a controlled setting using the executable query layer. Extracted strategies were normalized into a canonical Boolean form and executed against OpenAlex where feasible, with retrieved results compared to resolved reference lists. In this demonstration, normalized Boolean strategies achieved higher recall than keyword-based queries (0.245 vs.\ 0.180) at similarly low precision (0.015). These signals do not aim to reproduce database-specific search performance. They should be interpreted as retrieval behaviour within a simplified, uniform index that is suitable for cross-domain comparison and method development.
% Within this setting, extracted queries can be normalised, executed against OpenAlex where feasible, and evaluated against resolved reference sets without confounding effects from heterogeneous database infrastructures. In the demonstration retrieval experiment, normalized Boolean strategies achieved higher recall than keyword-based queries (0.245 vs.\ 0.180) at similarly low precision (0.015). These results do not aim to reproduce database-specific performance but reflect retrieval behaviour under a simplified, cross-domain index. More broadly, the corpus makes it possible to examine cross-disciplinary generalization, compare retrieval approaches under controlled conditions, and analyze how disciplinary vocabulary, logical structure, and normalization choices influence retrieval performance.

Beyond this retrieval demonstration, the corpus supports use cases that do not depend on executable queries. For screening research, the reference backbone and method fields support the study and benchmarking of prioritization and selection components in systematic review pipelines, while making explicit that reference lists are not identical to included-study sets. For extraction research, the structured fields derived from full text can serve as targets for models that identify and recover review artifacts such as objectives, eligibility criteria, and search strategies. For meta-research and science studies, Webis-SR4ALL-26 enables comparative analyses of systematic review practices across disciplines and time, and supports investigations into the often neglected role of systematic reviews in science communication and public policy~\cite{blumel:2020, mcmahan:2021, schniedermann:2022, simons:2021}.
% Beyond retrieval evaluation, the corpus also supports observational analysis of systematic review methodology at scale. The structured representation of objectives, eligibility criteria, and reported search strategies makes it possible to examine how reporting practices vary across domains and how search formulation conventions differ between review communities. Such analyses do not require executable queries and instead rely on the extracted methodological descriptions to characterize variation in evidence synthesis practices across fields.

Several limitations follow from the design choices underlying the corpus construction. The identification of systematic reviews relies on explicit title-based declarations, which favors precision but excludes reviews that do not self-identify using this terminology, which is a known challenge~\cite{moher:2007}. Full-text availability is incomplete, which constrains methodological extraction to a subset of records and contributes to uneven field coverage. In addition, the strict verification protocol applied during extraction prioritizes precision over recall, so methodological elements that are implicitly stated or distributed across multiple passages may remain unstructured. The resolved reference sets used for evaluation do not correspond directly to included-study gold standards, as reviews cite background and contextual literature in addition to eligible studies. Finally, the normalization of reported search strategies into simplified executable representations introduces abstractions that may not fully capture database-specific syntax, field restrictions, or manual refinements present in the original formulations.

Future work can address several of these constraints. Systematic review identification could be extended beyond title-based heuristics by incorporating additional metadata signals and full-text cues while maintaining comparable quality controls. Expanding full-text coverage would increase the proportion of reviews for which methodological information can be structured and analyzed. Further work is also needed to approximate or reconstruct included-study sets more directly, enabling evaluation settings that are closer to screening-level gold standards. Finally, improvements in extraction methods that increase recall without compromising evidence grounding would strengthen the completeness of the structured representation while preserving verifiability.

% The reported retrieval metrics quantify how well a canonical, executable approximation of a review’s reported search strategy can recover that review’s resolved reference set when run in OpenAlex. We treat the cited references as relevance signals because they represent studies that the review authors judged important enough to cite and discuss. At the same time, this set is not identical to the review’s included-study set, since systematic reviews also cite background and contextual literature, and they often identify additional eligible studies via citation chasing and other supplementary discovery practices that are not captured by database queries alone.

% The comparatively low precision and recall are therefore expected given the constraints of the evaluation setting. Our canonical queries are intentionally minimal and field-agnostic, and we remove database-specific syntax, field restrictions, and manual tuning to support consistent execution across domains and review styles. 
% Consequently, these metrics should be interpreted as indicators of retrieval reach under standardized execution conditions, not as a direct assessment of the effectiveness of the original search strategies or of systematic review quality.

% The corpus relies on strict title-based filtering that requires an explicit declaration of ``systematic review'' in the title. 
% This design choice favors high precision, but likely under-represents systematic reviews that do not self-identify that way, and future work could broaden identification beyond titles while maintaining quality controls.
